\section{Structure against operator: a factorial decomposition}
\label{sec:factorial}

The accuracy gap between a CNN--transformer hybrid and a flat state space
backbone is usually decomposed by quoting the hybrid's own ablation table:
so many points from the hierarchy, so many from the convolutional stem, and
so on. That decomposition was measured on one architecture. Applying it to
the gap between two architectures assumes that each component contributes the
same amount when the token-mixing operator underneath it is different, and
that assumption is precisely what a factorial design exists to test.

\subsection{Design}

We cross two factors. \emph{Structure} is flat versus hierarchical with
convolutional locality; \emph{operator} is self-attention versus a
bidirectional selective scan. The four cells are a flat attention model
(A), a flat state space model (B), a hierarchical attention model (C), and a
hierarchical state space model (D) that we built for this experiment.

D is C with the operator swapped and nothing else: the attention call inside
each block is replaced by a bidirectional Mamba mixer, while the local
perception unit and the inverted-residual feed-forward network stay as they
are, and the stem, patch embeddings, and classifier head are inherited
rather than reimplemented. This is what makes the design a factorial one---
$B \to D$ and $A \to C$ have to be the same manipulation, or the interaction
term measures something other than what its name says. For the same reason
we use a bidirectional scan rather than the four-directional scan of later
hierarchical SSM backbones: four directions would make $B \to D$ change the
number of scan directions as well as the structure.

D is a controlled reimplementation for the purpose of factor separation, not
a proposed architecture; hierarchical state space backbones already exist in
the literature.

All four cells are held to a parameter budget of \ParamBudget{}~M within
\ParamTolerance{} percent, and all four land inside it. Training uses one recipe for all four cells,
$300$ epochs, three seeds, on Tiny-ImageNet at $64 \times 64$.

\begin{figure}[tbp]
  \centering
  \includegraphics[width=\linewidth]{fig4_factorial.pdf}
  \caption{Cell means with standard deviation over three seeds, and the
  validation curves of all twelve runs. The two hierarchical cells separate
  from the two flat ones; the two operators do not separate from each other
  within either structure.}
  \label{fig:factorial}
\end{figure}

\begin{table}[t]
\centering
\caption{E4 cell means, Tiny-ImageNet top-1 (\%), $n=3$ seeds.}
\label{tab:e4cells}
\begin{tabular}{lrr}
\toprule
Cell & Top-1 & Std \\
\midrule
A DeiT-Ti (flat, attn) & 54.17 & 0.48 \\
B Vim-Ti (flat, SSM) & 52.57 & 0.94 \\
C CMT-Ti (hier, attn) & 60.19 & 0.48 \\
D H-Vim (hier, SSM) & 60.44 & 0.22 \\
\bottomrule
\end{tabular}
\end{table}

\begin{table}[tbp]
\centering
\caption{E4 factorial effects in percentage points, mean and standard deviation over $n=3$ seeds. Only the structure effect is large relative to its spread. The interaction sign is positive in every seed but its magnitude varies by a factor of 4.7, so we report the sign and not the size.}
\label{tab:e4effects}
\begin{tabular}{lrr}
\toprule
Effect & Mean & Std \\
\midrule
Structure $(C{+}D)/2-(A{+}B)/2$ & +6.94 & 0.64 \\
Operator $(A{+}C)/2-(B{+}D)/2$ & +0.68 & 0.45 \\
Interaction $(D{-}B)-(C{-}A)$ & +1.86 & 1.37 \\
\bottomrule
\end{tabular}
\end{table}


\subsection{What can be said}

The structure effect is \StructureEffect{} points with a standard deviation
of \StructureStd{} across seeds---about eleven times its own spread, and the
same size in every seed. This is the one quantity in the experiment that can
be quoted on its own. The direction agrees with the decomposition we set out
to check: most of the difference comes from the structural prior.

\subsection{What cannot}

The interaction is \InteractionEffect{} points with a standard deviation of
\InteractionStd{}. Its sign is positive in all three seeds, so the direction
survives; its magnitude varies between seeds by a factor of nearly five, so
the size does not. The correct statement is that the hierarchical structure
appears to help the state space operator at least as much as it helps
attention, in every seed we ran, by an amount this experiment does not
determine.

The spread has an identifiable source. The flat state space cell in the third
seed finished more than a point below its other two seeds, and since the
interaction is $(D-B)-(C-A)$, a low $B$ enters it directly. That cell has the
largest seed-to-seed standard deviation of the four.

The operator main effect is \OperatorEffect{} points with a standard
deviation of \OperatorStd{}---two thirds of the effect itself. We do not
quote its size either. What the experiment supports is weaker and still worth
saying: at this sequence length the choice of token-mixing operator matters
far less than the choice of structure.

\subsection{A registered prediction that failed, and what we did about it}

We registered a sanity band of 45 to 60 percent top-1 for all four cells
before running. Both hierarchical cells finished above it.

Exceeding a sanity bound from above is not self-evidently benign, so we
checked the pipeline rather than assuming. Two things rule out a defect.
All four cells draw from the same data loader, and both flat cells sit
inside the band---a label leak or a mapping error would lift all four
together. And the validation split itself was inspected directly: two
hundred classes, ten thousand images, exactly fifty per class, labels
covering the full range, and an evaluation transform with no augmentation,
reading from a directory disjoint from the training data.

The band was simply set too low for the hierarchical cells. We report the
prediction as registered and as missed; the pre-registration table was not
edited.

\subsection{An interaction that accuracy does not show and cost does}

Adding the hierarchy costs attention a little time and saves the state space
operator a great deal of it. Cell A trains in \HoursA{} hours and cell C in
\HoursC{}; cell B trains in \HoursB{} hours and cell D in \HoursD{}. The
attention model gets about seven percent slower, and the state space model
runs in a little over half the time.

The mechanism is the sequence length. In the hierarchical body the token
count halves at each stage, so the scan---whose cost is linear in sequence
length and carries a large constant---shortens stage by stage. Attention at
sixty-four tokens was already inexpensive, so the hierarchy buys it nothing
and the convolutional stem and local perception units cost it a little.

In accuracy the interaction between structure and operator is small and its
size undetermined. In wall-clock cost the interaction is unambiguous and
large. This was not in the pre-registration table, which concerned accuracy
only, and we report it as an observation made after the fact.

\subsection{Residual asymmetries}

Two differences between the cells are not eliminated by the parameter budget
and could in principle enter the effects.

The first is a normalisation and precision asymmetry that follows from how D
is built. Inheriting C's skeleton gives D's blocks layer normalisation with
residuals in half precision, whereas the flat state space cell B is
constructed as its authors intended, with RMS normalisation, a fused
add-norm, and residuals kept in single precision for numerical stability.
So $A \to C$ is layer norm to layer norm, while $B \to D$ is RMS norm with
fp32 residuals to layer norm with fp16 residuals. The two manipulations
differ in a way the factor names do not capture, and the difference lands on
the interaction term. It cannot be removed in code without D ceasing to share
C's skeleton, which would cost more than it buys; we record it rather than
claim it is absent.

The second concerns weight decay, and is described in Appendix~\ref{app:hvim}.
It falls mainly on the operator axis, and its cancellation is approximate
rather than exact.

Neither is large enough to explain the structure effect, which is an order of
magnitude above both. Both are the right size to matter for the interaction
term, which is a further reason we report its sign and not its size.
